\section{背景知识：x86 两级页表}

当 \reg{CR0} 的 PG 位（bit 31）被置 1 后，CPU 的每一次内存访问都会经过
\term{MMU}（Memory Management Unit）进行虚拟→物理地址翻译。
x86 32 位模式使用\term{两级页表}结构：

\begin{figure}[H]
  \centering
  % figures/ch08/fig01-vaddr-split.tex — auto-extracted from chapter source
\begin{tikzpicture}[
  box/.style={draw, minimum height=0.7cm, font=\ttfamily\small, anchor=west},
]
  % 虚拟地址 32-bit 分段
  \node[font=\bfseries\small] at (-1.5, 0.35) {虚拟地址};
  \node[box, minimum width=3cm, fill=red!15]   (pdi) at (0, 0) {PD index};
  \node[box, minimum width=3cm, fill=green!15] (pti) at (3, 0) {PT index};
  \node[box, minimum width=3cm, fill=blue!15]  (off) at (6, 0) {Offset};

  % bit 标注
  \node[font=\scriptsize, above] at (0, 0.35) {31};
  \node[font=\scriptsize, above] at (2.95, 0.35) {22};
  \node[font=\scriptsize, above] at (3.05, 0.35) {21};
  \node[font=\scriptsize, above] at (5.95, 0.35) {12};
  \node[font=\scriptsize, above] at (6.05, 0.35) {11};
  \node[font=\scriptsize, above] at (9, 0.35) {0};

  % 尺寸标注
  \node[font=\scriptsize\itshape, below] at (1.5, -0.05) {10 bits};
  \node[font=\scriptsize\itshape, below] at (4.5, -0.05) {10 bits};
  \node[font=\scriptsize\itshape, below] at (7.5, -0.05) {12 bits};
\end{tikzpicture}

  \caption{32 位虚拟地址的三段划分}
  \label{fig:vaddr-split}
\end{figure}

翻译流程如下：

\begin{enumerate}
  \item CPU 从 \reg{CR3} 取得 \term{Page Directory} 的物理基地址。
  \item 用虚拟地址的高 10 位（PD index）索引 Page Directory，得到一个 PDE。
  \item PDE 的高 20 位指向一个 \term{Page Table} 的物理地址。
  \item 用虚拟地址的中 10 位（PT index）索引 Page Table，得到一个 PTE。
  \item PTE 的高 20 位是目标\term{物理页帧}地址，拼上低 12 位 Offset，即为最终物理地址。
\end{enumerate}

\begin{figure}[H]
  \centering
  % figures/ch08/fig02-two-level-walk.tex — auto-extracted from chapter source
\begin{tikzpicture}[
  entry/.style={draw, minimum width=2.2cm, minimum height=0.55cm, font=\ttfamily\scriptsize},
  arr/.style={-{Stealth[length=3pt]}, thick},
  label/.style={font=\small\bfseries},
]
  % CR3
  \node[entry, fill=orange!20] (cr3) at (0, 3) {CR3};

  % Page Directory
  \node[label] at (3.5, 4.2) {Page Directory};
  \foreach \i/\c in {0/white, 1/white, 2/red!15, 3/white} {
    \node[entry, fill=\c] (pd\i) at (3.5, 3-\i*0.6) {PDE[\i]};
  }
  \node[font=\scriptsize] at (3.5, 0.4) {$\vdots$};
  \node[entry] (pd1023) at (3.5, -0.1) {PDE[1023]};

  % Page Table
  \node[label] at (7.5, 4.2) {Page Table};
  \foreach \i/\c in {0/white, 1/green!15, 2/white, 3/white} {
    \node[entry, fill=\c] (pt\i) at (7.5, 3-\i*0.6) {PTE[\i]};
  }
  \node[font=\scriptsize] at (7.5, 0.4) {$\vdots$};
  \node[entry] (pt1023) at (7.5, -0.1) {PTE[1023]};

  % Physical Page
  \node[entry, fill=blue!15, minimum width=2cm] (phys) at (11, 1.8) {物理页};

  % arrows
  \draw[arr] (cr3) -- (pd0);
  \draw[arr, red!60!black] (pd2) -- (pt0) node[midway, above, font=\scriptsize] {高 20 位};
  \draw[arr, green!50!black] (pt1) -- (phys) node[midway, above, font=\scriptsize] {帧地址};

  % index labels
  \node[font=\scriptsize\itshape, red!60!black, above=0.1cm of pd2, xshift=-1.2cm] {PD idx};
  \node[font=\scriptsize\itshape, green!50!black, above=0.1cm of pt1, xshift=-1.2cm] {PT idx};
\end{tikzpicture}

  \caption{两级页表翻译流程：CR3 → PDE → PTE → 物理页}
  \label{fig:two-level-walk}
\end{figure}

\begin{whybox}
为什么用两级而不是一级平坦页表？

一级页表需要 $2^{20}$ 个 PTE = 4\,MiB 连续内存，即使地址空间大部分为空也要全部分配。
两级结构只需 1 个 Page Directory（4\,KiB），空闲区域的 PDE 标记为"不存在"，
对应的 Page Table 根本不需要分配。大多数进程只用几十个 Page Table，远小于 4\,MiB。
\end{whybox}

% ============================================================================

